\documentclass[12pt, a4paper]{scrartcl}
\usepackage{azzam}

\begin{document}

\section*{Mini Simulation 2 : OSK SMP 2026 6 Mei}
\begin{enumerate}
    \item Diketahui $a$ dan $b$ dua bilangan real nonnegatif yang memenuhi $a+2b=15$. Nilai terbesar yang mungkin dari $3a+5b$ adalah . . . .
    \begin{enumerate}[label={(\Alph*)}]
        \item 49 \item 45 \item 51 \item 41
    \end{enumerate}

    \item Dua benteng berbeda diletakkan pada sembarang petak di papan catur $8\times 8$. Dua buah benteng dikatakan \textit{saling menyerang} jika berada pada baris atau kolom yang sama. Peluang kedua benteng tidak saling menyerang adalah . . . .
    \begin{enumerate}[label={(\Alph*)}]
        \item $\dfrac{49}{64}$ \item $\dfrac{7}{9}$ \item $\dfrac{3}{4}$ \item $\dfrac{16}{21}$
    \end{enumerate}

    \item Diberikan persegi panjang $ABCD$ dengan luas $300$ satuan luas. Titik $E$ dan $F$ berturut-turut terletak pada sisi $AB$ dan $BC$ sedemikian sehingga berlaku perbandingan panjang $AE:EB=1:2$ dan $BF:FC=1:1$. Luas dari segiempat $DEFC$ adalah . . . satuan luas.
    \begin{enumerate}[label={(\Alph*)}]
        \item 150 \item 200 \item 160 \item 240
    \end{enumerate}

    \item Suatu bilangan asli disebut $bagus$ jika dapat dinyatakan sebagai $5a+6$, $6b+13$, dan $7c+1$ sekaligus untuk suatu bilangan bulat $a$, $b$, dan $c$. Perhatikan bahwa $1$ merupakan bilangan bagus karena
    $$1 = 5\times (-1)+6 = 6\times (-2) +13 = 7\times 0+1.$$
    Misalkan $N$ adalah bilangan asli terkecil selanjutnya yang bagus. Jumlah digit-digit dari $N$ adalah . . . .
    \begin{enumerate}[label={(\Alph*)}]
        \item 3 \item 4 \item 5 \item 6
    \end{enumerate}

    \item Dalam lima hari, koperasi $A$ berhasil menjual sebanyak $2,3,5,8$, dan $12$ seragam sekolah secara berturut-turut. Ketua koperasi ingin mengambil $2$ bilangan dari kelima bilangan tersebut, sehingga rata-rata dari dua bilangan yang dipilih lebih besar daripada rata-rata penjualan koperasi tersebut selama $5$ hari. Banyaknya cara yang mungkin dilakukan olehnya adalah . . . .
    \begin{enumerate}[label={(\Alph*)}]
        \item 2 \item 3 \item 4 \item 5
    \end{enumerate}

    \item Diketahui bahwa $x^2+px+q=0$ mempunyai $2$ akar bulat positif. Jika $p+q=4,$ maka nilai dari $q$ adalah . . . .
    \begin{enumerate}[label={(\Alph*)}]
        \item 8 \item 10 \item 12 \item 14
    \end{enumerate}

    \item Jin-ho mempunyai $64$ kubus berukuran $1\times 1\times 1$. Dari sekumpulan kubusnya, Jin-ho hendak menyusun seluruh kubus miliknya sehingga menjadi sebuah balok. Dua balok dikatakan berbeda jika salah satu balok tidak dapat dihasilkan dari refleksi, rotasi, ataupun translasi balok yang lainnya. Banyaknya macam balok berbeda yang bisa disusun oleh Jin-ho adalah . . . .
    \begin{enumerate}[label={(\Alph*)}]
        \item 4 \item 5 \item 6 \item 7
    \end{enumerate}

    \newpage
    \item Diberikan trapesium sama kaki $ABCD$ dengan $AB$ sejajar $CD$. Diketahui $\angle ACB=\angle ADB=30^\circ$ dan panjang $DA=AB=BC=1$. Luas dari $ABCD$ adalah . . . .
    \begin{enumerate}[label={(\Alph*)}]
        \item $\sqrt{3}$ \item $\dfrac{\sqrt{3}}{2}$ \item $\dfrac{3\sqrt{3}}{4}$ \item $\dfrac{\sqrt{3}}{4}$
    \end{enumerate}

    \item Banyaknya pasangan bilangan prima terurut $(p,q)$ sehingga $24-p-q$ juga bilangan prima adalah . . . .
    \begin{enumerate}[label={(\Alph*)}]
        \item 11 \item 13 \item 15 \item 17
    \end{enumerate}

    \item Andi memiliki $4$ koin yang memiliki $2$ sisi: sisi angka dan sisi gambar. Andi ingin mengetos keempat koin tersebut. Diketahui bahwa jika satu koin ditos, peluang munculnya sisi angka adalah $\frac{2}{5}$. Peluang bahwa agar terdapat minimal dua koin yang menunjukkan sisi gambar dalam pengetosan keempat koin tersebut adalah . . . .
    \begin{enumerate}[label={(\Alph*)}]
        \item $\dfrac{513}{625}$ \item $\dfrac{112}{625}$ \item $\dfrac{171}{625}$ \item $\dfrac{441}{625}$
    \end{enumerate}
\end{enumerate}

\end{document}